\documentclass{article}
\usepackage{graphicx} % Required for inserting images

\title{Math}
\author{Andres Soto Ticse}
\date{August 2025}

\begin{document}

\maketitle

\section{Introduction}

\paragraph{Teorema FWL}
\begin{itemize}
  \item Sea el modelo lineal
  \begin{equation}
    y = X_1\beta_1 + X_2\beta_2 + u. \label{eq:modelo}
  \end{equation}
  con $[X_1\ X_2]$ de rango completo.

  \item Matriz de proyección sobre $X_2$:
  \begin{equation}
    P_2 = X_2(X_2'X_2)^{-1}X_2'. \label{eq:P2}
  \end{equation}

  \item Matriz de aniquilación de $X_2$:
  \begin{equation}
    M_2 = I - P_2. \label{eq:M2}
  \end{equation}

  \item Propiedades: $P_2 = P_2' = P_2^2$, \ $M_2 = M_2' = M_2^2$, \ $M_2X_2 = 0$.
\end{itemize}

Las ecuaciones normales del MCO sobre $[X_1\ X_2]$ son
\begin{align}
\begin{equation}
X_1'(y - X_1\beta_1 - X_2\beta_2) &= 0, \label{eq:normal1}
\end{equation}

\begin{equation}
 X_2'(y - X_1\beta_1 - X_2\beta_2) &= 0. \label{eq:normal2}
\end{equation}
 
\end{align}


Apartir de (5):
\begin{equation}
\beta_2 = (X_2'X_2)^{-1}X_2'(y - X_1\beta_1). \label{eq:beta2}
\end{equation}
\indent Sustituyendo en (4) y usando M2 se obtiene
\begin{equation}
X_1'M_2\,y = X_1'M_2X_1\,\beta_1. \label{eq:ecuacionM2}
\end{equation}

Dado que $X_1'M_2X_1$ es no singular, el estimador es
\begin{equation}
\hat\beta_1=(X_1'M_2X_1)^{-1}X_1'M_2y. \label{eq:beta1hat}
\end{equation}

Definiendo $\tilde y=M_2y$ y $\tilde X_1=M_2X_1$, las ecuaciones normales de la regresión 
\indent MCO de $\tilde y$ sobre $\tilde X_1$ son


\begin{equation}
\tilde X_1'(\tilde y-\tilde X_1\beta_1)=0
\ \Longleftrightarrow\
X_1'M_2y=X_1'M_2X_1\,\beta_1, \label{eq:tilde-normal}
\end{equation}
\indent que coinciden con M2. Por tanto, el coeficiente al regredir $\tilde y$ sobre $\tilde X_1$ es 
\indent exactamente $\hat\beta_1$ de la regresión completa.



\end{document}
